\section{Introduction}

The convergence of large language models with high-performance computing is transforming scientific workflows. Autonomous agents are beginning to appear alongside established workflow automation systems for data-intensive science, orchestration, and large-scale experiment management~\cite{deelman2015pegasus,babuji2019parsl,jain2015fireworks,xi2023agentsurvey}. These AI agents promise to accelerate discovery, but they also introduce a threat category that HPC security was never designed to address.

The core problem is not that AI agents are untrustworthy. The problem is that they are too trustworthy, specifically, they trust the wrong inputs. Recent work has shown that LLM-integrated applications and tool-using agents are vulnerable to prompt injection and instruction-channel confusion~\cite{greshake2023indirect,liu2024prompt,liu2023houyi}. An agent authorized to process genomics data for Project X will faithfully execute whatever instructions it receives, including instructions embedded within that data by an attacker. A prompt injection payload hidden in a \textit{FITS} file header, a malicious command left in a shared \textit{/tmp} directory by a co-located job, or a compromised tool that returns hidden instructions during normal operation represent concrete attack vectors. These are not theoretical concerns; they are the natural consequence of deploying instruction-following systems into environments where untrusted data is routine.

This paper identifies and formalizes a threat that, to our knowledge, has not been studied in prior work for HPC specifically: \textit{the hijacked authorized agent}. Unlike rogue agents, which authentication blocks, or adversarial agents, which authorization constrains, a hijacked agent operates with the full credentials and privileges of a legitimate user. It does not need to escalate privileges because it already has them. It does not need to bypass access controls because they already permit everything the user can do. From the perspective of every existing monitoring system that validates identity, the agent is legitimate. Yet it exfiltrates data through precisely the channel no HPC system inspects: the encrypted HTTPS connection to its LLM backend.

The attack surface is uniquely HPC. Shared parallel filesystems such as Lustre and IBM Spectrum Scale create injection vectors with no direct web analogue~\cite{lustre2025intro,ibm2024spectrumscale}. Multi-tenant compute nodes scheduled by \slurm{} create co-location injection opportunities through shared temporary storage and kernel resources~\cite{yoo2003slurm}. Tool-use protocols and agent integration layers make it easier to connect agents to external services and local executables, but they also widen the trusted computing base~\cite{mcp2024spec,mialon2023augmented}. Coordinated multi-agent workflows enable distributed exfiltration that evades per-agent reasoning entirely.

Existing defenses are fundamentally ill-equipped against this threat. Authentication succeeds because the agent has valid Kerberos tickets. Authorization succeeds because the agent has legitimate POSIX permissions. Network monitoring is blinded because LLM API traffic is encrypted. Data loss prevention fails because sensitive data can be embedded within whitelisted HTTPS requests. User behavior analytics report nothing because the agent's workflow patterns remain superficially consistent with its authorized role. The defender's toolkit checks identity, but not intent.

We propose \textbf{behavioral attestation} as the foundation for securing AI agents in HPC. Our key insight is straightforward: we cannot prevent all injection attacks, because the attack surface is too large and diverse. However, we can contain their effects by detecting and responding to the behavioral violations that hijacked agents inevitably produce. When an agent accesses files outside its authorized project, connects to non-whitelisted endpoints, invokes unauthorized tools, or exceeds its data budget, attestation detects the violation and enforces containment before larger-scale exfiltration can continue unchecked.

Behavioral attestation differs from prior system-integrity attestation approaches along three fundamental dimensions. First, it provides constraint-based verification rather than probabilistic alerts. Existing agent defenses often rely on classifiers or pattern-matching systems with inherent false positive and false negative rates. Behavioral attestation asks whether an action satisfies the declared policy boundary. Second, it is constraint-based rather than signature-based. Existing defenses often block known-bad patterns, which attackers can evade through obfuscation. Constraints define what is allowed, not what is malicious. Third, it operates at runtime rather than post hoc. Existing monitoring systems often alert after the attack succeeds; behavioral attestation verifies constraints continuously during execution and triggers containment while the session is still active.

This paper makes the following contributions.

\begin{itemize}
    \item We formalize the hijacked authorized agent threat in HPC, identifying four distinct attack vectors that exploit unique characteristics of HPC infrastructure.
    \item We introduce behavioral attestation as a constraint-based runtime security primitive that extends zero-trust principles from identity to agent behavior.
    \item We design and implement AEGIS, a complete zero-trust architecture encompassing constraint specification, \ebpf-based attestation, policy verification, cluster-wide correlation for coordinated attack detection, and automated containment through \slurm{} integration.
    \item We demonstrate that \projectname~detects all four implemented attack vectors on the real attestation path, achieves full coverage in the controlled attack suite, outperforms representative baselines, and maintains low measured runtime overhead on \epyc{} hardware.
\end{itemize}
